\documentclass[11pt]{article}
\usepackage[a4paper,margin=22mm]{geometry}
\usepackage{amsmath,amssymb,mathtools,bm}
\usepackage{booktabs,longtable,array,microtype,xurl,hyperref}
\hypersetup{colorlinks=true,linkcolor=blue,citecolor=blue,urlcolor=blue,
 pdftitle={P54PQ-v2 P1 stationary points, Hessian, and spectrum}}
\newcommand{\Spin}{\operatorname{Spin}}
\newcommand{\tr}{\operatorname{tr}}
\newcommand{\cH}{\mathcal H}
\newcommand{\cM}{\mathcal M}
\newcommand{\hc}{\mathrm{h.c.}}
\title{P1: Stationary Points, Goldstone Alignment,\\
Complete Hessian, and Spectrum in \texttt{P54PQ-v2}}
\author{Route-F four-dimensional mainline}
\date{30 August 2026}

\begin{document}
\maketitle

\begin{abstract}
We construct P1 directly from the corrected \texttt{P54PQ-v2} tensor
action.  The differentiable field space has
$54+2(126)+2(10)+2=328$ canonically normalized real coordinates.  We derive
the radial stationary equations in the action card's own normalization,
enumerate competing radial branches, construct the complete
$328\times328$ scalar Hessian by automatic differentiation, and independently
construct the 45-generator gauge orbit.  Gauge invariance requires 33 gauge
zero modes; the faithful PQ symmetry adds one physical axion direction after
projecting away the gauge orbit.  The calculation is fail-closed: a
stationary saddle is reported as a saddle, and the electroweak doublet tuning
is not silently identified with vacuum stability.
\end{abstract}

\section{Why P1 is restarted on version 2}

The independent invariant audit found
\begin{equation}
 V_7=\chi_7\Phi_{ij}\phi_i\phi_jS^*+\chi_7^*\Phi_{ij}\phi_i^*\phi_j^*S.
 \label{eq:chi7}
\end{equation}
It contributes at second order in the $10_H$ fluctuation whenever
$\langle\Phi\rangle$ and $\langle S\rangle$ are nonzero.  Therefore it changes
the $10_H$ diagonal and holomorphic mass blocks even though it does not alter
the high-scale radial stationarity equations at $\langle\phi\rangle=0$.
Computing a ``complete'' Hessian from version~1 and appending \eqref{eq:chi7}
after diagonalization would be mathematically invalid.

The earlier literature potential also contains printed conjugation and index
errors, and its displayed radial coefficients are not all reproduced by its
displayed factorials under a single tensor normalization.  P1 consequently
uses one source of truth: the explicit tensor action in the corrected action
card.  Literature mass tables are comparisons, not definitions.

\section{Canonical field coordinates}

Let $E^a_{ij}$, $a=1,\ldots,54$, be an orthonormal basis of real symmetric
traceless matrices,
\begin{equation}
 E^a_{ij}E^b_{ij}=\delta^{ab},\qquad \Phi_{ij}=x_aE^a_{ij}.
\end{equation}
Let $U_{A\alpha}$ be an orthonormal basis of one Hodge eigenspace in the 252
independent five-form components, $\alpha=1,\ldots,126$.  The complex fields
are written
\begin{align}
 z_\alpha&=\frac{x_{54+\alpha}+i x_{180+\alpha}}{\sqrt2},
 &\Sigma_A&=U_{A\alpha}z_\alpha,\nonumber\\
 h_i&=\frac{x_{306+i}+i x_{316+i}}{\sqrt2},
 &S&=\frac{x_{326}+ix_{327}}{\sqrt2}.
\end{align}
In these variables the kinetic term is $\frac12\partial x_I\partial x_I$.
The full scalar coordinate count is
\begin{equation}
 \boxed{N_{\mathbb R}=54+252+20+2=328.}
\end{equation}

The high-scale background is
\begin{align}
 \Phi_0&=\omega X,\qquad
 X=\operatorname{diag}(-2/5,-2/5,-2/5,-2/5,-2/5,-2/5,
 3/5,3/5,3/5,3/5),\nonumber\\
 \Sigma_0&=\sigma\,\Omega,qquad
 \|\Omega\|^2_{\wedge^5}=1,qquad
 S_0=\frac{v_s}{\sqrt2},\qquad h_0=0.
 \label{eq:background}
\end{align}
The ordered component of the holomorphic form has magnitude
$|\Omega_{246810}|=1/(4\sqrt2)$, matching the action card.

\section{Stationary equations in the card normalization}

Direct substitution of \eqref{eq:background} into the tensor action gives
\begin{align}
V_{\rm rad}={}&-\frac65\mu^2\omega^2+\frac4{25}c\omega^3
+A_4\omega^4-\frac12\nu^2\sigma^2+\frac14\lambda_0\sigma^4
+\frac65d\,\omega^2\sigma^2\nonumber\\
&-\frac12\mu_s^2v_s^2+\frac14\chi_1v_s^4
+60\chi_2\sigma^2v_s^2+\frac65\chi_3\omega^2v_s^2,
\label{eq:vrad}\
A_4={}&\frac{36}{25}a+\frac{42}{125}b,qquad d=\alpha-\beta.
\end{align}
The $\lambda_{2,4,4'}$ and $\chi_4$ contractions vanish on the decomposable
holomorphic five-form, but their second variations do not vanish and they
remain in the full Hessian.  Equation~\eqref{eq:vrad} is not imported from a
radial formula; it is an executable contraction fingerprint of the card.

For the fully nonzero branch, stationarity is equivalent to
\begin{align}
\mu^2&=\frac15c\omega+\frac53A_4\omega^2+d\sigma^2+\chi_3v_s^2,
\label{eq:stat1}\\
\nu^2&=\lambda_0\sigma^2+\frac{12}{5}d\omega^2
+120\chi_2v_s^2,
\label{eq:stat2}\\
\mu_s^2&=\chi_1v_s^2+120\chi_2\sigma^2
+\frac{12}{5}\chi_3\omega^2.
\label{eq:stat3}
\end{align}
The factors of 60 and 120 are not physical enhancements; they are the
declared normalization of the coefficient $\chi_2$ in the tensor card.  A
rescaled coefficient $\widehat\chi_2=120\chi_2$ would remove them, but mixing
the two parameter conventions is forbidden.

The radial curvature before conversion to the canonical metric is
\begin{equation}
K_{\rm rad}=\begin{pmatrix}
\frac{12}{25}c\omega+8A_4\omega^2&\frac{24}{5}d\omega\sigma&\frac{24}{5}\chi_3\omega v_s\\
\frac{24}{5}d\omega\sigma&2\lambda_0\sigma^2&240\chi_2\sigma v_s\\
\frac{24}{5}\chi_3\omega v_s&240\chi_2\sigma v_s&2\chi_1v_s^2
\end{pmatrix}.
\label{eq:radialH}
\end{equation}
Its generalized eigenvalues use
$G_{\rm rad}=\operatorname{diag}(12/5,2,1)$, since the coordinates
$(\omega,\sigma,v_s)$ are not all canonical.  Positivity of
\eqref{eq:radialH} is necessary but not sufficient for a vacuum.

Every subset of $\{\omega,\sigma,v_s\}$ can vanish, so the radial problem has
eight algebraic branch types before sign and multiplicity are resolved.  The
machine audit starts from a grid of positive, negative, and zero initial
conditions, deduplicates all roots, evaluates $V_{\rm rad}$, and records the
three radial curvatures.  This prevents the desired nonzero branch from being
declared the vacuum merely because its mass squares were solved from
\eqref{eq:stat1}--\eqref{eq:stat3}.

\section{Complete Hessian}

The scalar Hessian is defined without a representation-level truncation:
\begin{equation}
 \boxed{\cH_{IJ}=\left.\frac{\partial^2V}{\partial x_I\partial x_J}
 \right|_{x=x_0},\qquad I,J=1,\ldots,328.}
 \label{eq:fullH}
\end{equation}
Automatic differentiation is applied to every indexed contraction in
$V_{\Phi\Sigma}+V_{\Phi\Sigma\phi}+V_S$, including all Hermitian partners.
The block layout is
\begin{equation}
\cH=\begin{pmatrix}
H_{54,54}&H_{54,126}&H_{54,10}&H_{54,S}\\
H_{126,54}&H_{126,126}&H_{126,10}&H_{126,S}\\
H_{10,54}&H_{10,126}&H_{10,10}&H_{10,S}\\
H_{S,54}&H_{S,126}&H_{S,10}&H_{S,S}
\end{pmatrix}_{(54+252+20+2)\times(54+252+20+2)}.
\end{equation}
At $h_0=0$, symmetry makes some displayed blocks vanish, but they are
computed rather than omitted.  In particular, $\eta_1$ can mix the 126 and
10 sectors at quadratic order, $\chi_4$ mixes singlet sectors, and
\eqref{eq:chi7} changes the $10$ block.  The verifier repeats the Hessian with
$\chi_7=0$ and requires a nonzero Frobenius difference.

\section{Goldstone alignment theorem and its numerical form}

For a generator $T_A$, gauge invariance says
\begin{equation}
0=\frac{d}{dt}V(e^{tT_A}x)\big|_{t=0}
=\partial_IV(x)(T_Ax)^I.
\end{equation}
Differentiating with respect to $x^J$ and imposing stationarity gives
\begin{equation}
 \cH_{JI}(T_Ax_0)^I=0.
 \label{eq:goldstoneproof}
\end{equation}
Thus the theorem compares actual vectors, not just the number of small
eigenvalues.  The 45-column orbit matrix is
\begin{equation}
 \cM_G=\bigl(T_1x_0\ \cdots\ T_{45}x_0\bigr).
\end{equation}
The independent tensor construction gives
\begin{equation}
 \operatorname{rank}\cM_G(\Phi_0)=24,qquad
 \operatorname{rank}\cM_G(\Phi_0,\Sigma_0)=33,
\end{equation}
so the unbroken algebra has dimension $45-33=12$, as required for
$SU(3)_C\times SU(2)_L\times U(1)_Y$.

Let $q x_0$ be the infinitesimal PQ direction.  Its component parallel to
the gauge orbit is unphysical, so the physical axion direction is
\begin{equation}
 a_0=\left[1-\cM_G\cM_G^+\right]q x_0,
 \label{eq:axionproject}
\end{equation}
where $+$ denotes the Moore--Penrose inverse.  P1 requires
\begin{equation}
 \operatorname{rank}(\cM_G,a_0)=34,qquad
 \frac{\|\cH\cM_G\|}{\|\cH\|\|\cM_G\|}\ll1,qquad
 \frac{\|\cH a_0\|}{\|\cH\|\|a_0\|}\ll1.
\end{equation}
The QCD anomaly does not enter this tree-level Hessian; it lifts the axion
only in the low-energy nonperturbative effective action.

\section{Scalar and vector spectra}

The scalar spectrum is the complete eigenvalue multiset
\begin{equation}
 \operatorname{Spec}_{\rm scalar}=\operatorname{eigvalsh}(\cH).
\end{equation}
Degeneracies are grouped only after the Goldstone residual fixes an absolute
tolerance.  A negative physical eigenvalue classifies the background as a
stationary saddle even when the radial $3\times3$ block is positive.

The vector spectrum follows from the same orbit matrix and not from the
potential:
\begin{equation}
 \frac{M_V^2}{g_{10}^2}=\cM_G^T\cM_G.
 \label{eq:vectorM}
\end{equation}
Its nullity must be 12 and its positive rank 33.  Equation~\eqref{eq:vectorM}
also fixes the generator normalization and prevents an arbitrary rescaling
of the scalar Goldstone vectors.

This P1 benchmark is dimensionless and exercises every Hessian block; it is
not a flavor, unification, or proton-decay fit.  It has no negative physical
scalar mode and exactly one tuned complex electroweak doublet.  The doublet
tuning is a rank condition on the full mixed doublet block, not an
independent free mass.

\section{Executed benchmark and complete spectrum}

We measure all dimensionful quantities in units of $\omega=1$ and choose
\begin{align*}
&(\sigma,v_s)=(0.35,0.25),\quad
(a,b,c)=(0.50,0.40,-0.10),\\
&(\lambda_0,\lambda_2,\lambda_4,\lambda'_4)=(0.90,0.80,0.60,0.40),
\quad(\alpha,\beta)=(0.25,0.10),\\
&(\chi_1,\chi_2,\chi_3,\chi_4,\chi_5,\chi_6,\chi_7)
=(0.80,0.001,0.05,0.03,0.20,0.025,0.03),\\
&(\eta_0,\eta_1,\eta_2,\eta_3)=(0.20,0.025,0.10,0.018),
\quad(\gamma_1,\gamma_2)=(0.12,-0.07),\\
&(\xi_1,\xi_2,\xi_3)=(0.35,0.08,-0.20).
\end{align*}
Equations~\eqref{eq:stat1}--\eqref{eq:stat3} then give
\begin{equation}
 (\mu^2,\nu^2,\mu_s^2)=(1.4255,0.47775,0.1847).
\end{equation}
The single doublet determinant condition fixes
\begin{equation}
 \xi_0^2=0.14616372381894022.
\end{equation}
This decimal is an exported root of the matrix rank condition, not an input
used to fit a measured mass.  At the present high-scale tree-level gate the
light eigenvalue is exactly zero; P2--P3 will replace that condition by the
renormalized electroweak matching condition.

The full gradient residual is $3.80\times10^{-16}$.  The radial solver finds
27 signed stationary roots; the desired $\omega=+1$ branch has
$V=-0.87681703125$ and is the lowest enumerated branch.  The Hessian has 38
zero modes and no negative modes at tolerance $8.50\times10^{-7}$.  Direct
projection accounts for 33 gauge modes, one PQ mode, and four light-doublet
real modes.  For the latter,
\begin{equation}
 (C_3,C_2,Y^2)=(0,3/4,1/4),qquad
 \|h_{10}\|^2=0.9999920937,quad
 \|h_{126}\|^2=7.9063\times10^{-6}.
\end{equation}
Thus it is one $(1,2,\pm1/2)$ real multiplet, not an accidental quartet.

The complete Standard-Model Casimir decomposition is below.  ``3-pair'' and
``6-pair'' denote an irrep and its conjugate in the real Hessian.  Equal
masses on distinct rows are genuine degeneracies, not merged labels.
\begin{longtable}{@{}rrrrr@{}}
\toprule
$m^2$ & $SU(3)$ & $\dim SU(2)$ & $|Y|$ & real mult.\\
\midrule
\endhead
$0$&1&1&0&2\\
$0$&1&2&$1/2$&4\\
$0$&3-pair&2&$1/6$&12\\
$0$&3-pair&1&$2/3$&6\\
$0$&1&1&1&2\\
$0$&3-pair&2&$5/6$&12\\
$0.0382803065$&1&2&$1/2$&4\\
$0.0888808155$&1&1&0&1\\
$0.1121882733$&3-pair&2&$1/6$&12\\
$0.1159441884$&1&1&0&1\\
$0.1866355208$&3-pair&1&$1/3$&6\\
$0.2116296538$&3-pair&1&$1/3$&6\\
$0.2600000000$&8&1&0&8\\
$0.2722422827$&6-pair&1&$2/3$&12\\
$0.3939174318$&1&2&$1/2$&4\\
$0.4006848071$&3-pair&1&$1/3$&6\\
$0.5400000000$&1&3&0&3\\
$0.5522408476$&1&3&1&6\\
$1.4210000000$&6-pair&1&$1/3$&12\\
$1.4210000000$&3-pair&3&$1/3$&18\\
$1.4210000000$&3-pair&1&$4/3$&6\\
$1.4700000000$&6-pair&3&$1/3$&36\\
$1.4700000000$&6-pair&1&$4/3$&12\\
$1.4700000000$&1&1&2&2\\
$1.5188026219$&8&2&$1/2$&32\\
$1.5188026219$&3-pair&2&$7/6$&12\\
$1.5210998141$&1&2&$1/2$&4\\
$1.5721973781$&8&2&$1/2$&32\\
$1.5721973781$&3-pair&2&$7/6$&12\\
$1.7576694975$&3-pair&1&$1/3$&6\\
$1.8266780732$&3-pair&1&$1/3$&6\\
$2.0580077173$&6-pair&1&$2/3$&12\\
$2.0580091524$&1&3&1&6\\
$2.1580617267$&3-pair&2&$1/6$&12\\
$2.8334249961$&1&1&0&1\\
\bottomrule
\end{longtable}
The real multiplicities sum to 328, and the maximum leakage of any mass
eigenspace under the 12 unbroken generators is $1.93\times10^{-10}$.
The independent vector result is
\begin{equation}
\frac{m_V^2}{g_{10}^2}:quad
0^{\times12},\quad0.1225^{\times8},\quad0.6125^{\times1},
\quad1^{\times12},\quad1.1225^{\times12}.
\end{equation}
The $\chi_7$-only Hessian has Frobenius norm $2.32379\times10^{-2}$, which
explicitly verifies that the invariant-basis correction enters P1.

\section{Acceptance and next gate}

P1 has four logically separate outputs:
\begin{enumerate}
\item a complete stationary-branch ledger from the corrected action;
\item a $328\times328$ Hessian whose field count and symmetry are checked;
\item vector-by-vector Goldstone alignment for 33 gauge modes and one
projected PQ mode;
\item the full scalar and vector eigenvalue multisets for the declared point.
\end{enumerate}
The executed benchmark passes all four items: the verifier reports 14/14
checks, a semidefinite radial-global branch, and exactly the required light
doublet.  This closes the tree-level algebraic P1 gate.  It does not assert
loop stability, a measured-scale fit, or cosmological longevity outside the
enumerated singlet branches.  The next mainline task is P2: run and match the
actual spectrum, propagate parameter/scale covariance, and reimpose the
light-doublet condition on the renormalized mass matrix.

\begin{thebibliography}{9}
\bibitem{BabuKhan2015}
K.~S.~Babu and S.~Khan,
``A Minimal Non-Supersymmetric SO(10) Model,''
\href{https://arxiv.org/abs/1507.06712}{arXiv:1507.06712}.
\end{thebibliography}
\end{document}
